% ============================================================================
% 明镜高悬 (MingJing)：基于 eBPF 的自主编程智能体跨层行为归因、
% 结果感知语义检测与多层操作系统强制访问控制
%
% 适用投递目标：《信息安全学报》《软件学报》《计算机研究与发展》及竞赛学术论文附件
% 编译方式：xelatex mingjing_paper_zh.tex
% ============================================================================
\documentclass[UTF8,a4paper,zihao=-4,twocolumn]{ctexart}

\usepackage[top=2.2cm,bottom=2.2cm,left=1.9cm,right=1.9cm,columnsep=0.65cm]{geometry}
\usepackage{amsmath,amssymb}
\usepackage{booktabs,tabularx,multirow,array}
\usepackage{enumitem}
\usepackage{url}
\usepackage[hidelinks]{hyperref}
\usepackage{fancyhdr}
\usepackage{caption}

\captionsetup{font=small,labelfont=bf}
\setlist{nosep,leftmargin=1.5em}
\setlength{\parskip}{2pt}
\pagestyle{fancy}
\fancyhf{}
\renewcommand{\headrulewidth}{0.4pt}
\fancyhead[L]{\zihao{-5}\songti 明镜高悬（MingJing）：基于 eBPF 的智能体跨层行为归因与多层安全控制}
\fancyhead[R]{\zihao{-5}\songti 学术论文稿（2026年11月修订版）}
\fancyfoot[C]{\zihao{5}\thepage}

\ctexset{
  section={format=\heiti\zihao{4},number=\arabic{section}},
  subsection={format=\heiti\bfseries\zihao{5},number=\thesection.\arabic{subsection}},
  subsubsection={format=\heiti\bfseries\zihao{-5},number=\thesubsection.\arabic{subsubsection}}
}

\newcommand{\code}[1]{\texttt{\small #1}}

\title{\heiti\zihao{3} 明镜高悬：基于 eBPF 的自主编程智能体跨层行为归因、\\结果感知语义检测与多层内核强制控制方法}

\author{\songti\zihao{-4} 匿名作者 / 项目研究组\\
\zihao{5}\textit{（投稿《信息安全学报》/《软件学报》专刊及双盲审稿预备稿）}}
\date{\zihao{5}2026年11月19日（基于仓库实测数据与安全审计更新）}

\begin{document}
\twocolumn[
\begin{@twocolumnfalse}
\maketitle
\begin{abstract}
\noindent
基于大语言模型（LLM）的自主编程智能体（如 Claude Code、Codex、Gemini CLI）与模型上下文协议（MCP）工具链能够将自然语言指令转化为跨解释器、Shell、构建工具与外部网络的复杂操作系统动作。然而，现有防护体系存在显著的\textbf{“语义—执行鸿沟”（Semantic-Execution Gap）}：应用层护栏仅检查工具调用声明或提示词，无法观测脱离运行时框架的派生子进程；而传统 eBPF 内核监控工具（如 Falco、Tetragon）仅捕获系统调用事实，缺乏智能体运行（\code{agent\_run\_id}）与工具调用（\code{tool\_call\_id}）的语义边界，且在多智能体并发调用同名程序时极易产生误归因。

针对上述挑战，本文设计、实现并评估了面向 Linux 研发环境的智能体行为观测与多层安全控制系统——\textbf{明镜高悬（MingJing / Agent eBPF Filter）}。首先，提出\textbf{三模态显式上下文接入与有界归因架构}，融合 Python/Node.js 进程注册适配器、10 类主流智能体 CLI 原生生命周期钩子以及基于 \code{SO\_PEERCRED} 鉴权的 UDS 命令包装器（\code{agent-wrapper}），在默认关闭高风险 TLS 明文劫持的前提下实现内核事件的跨层归因；其次，构建\textbf{面向取证完整性的遥测与结果感知语义检测管线}，提出早于 \code{bpf\_ringbuf\_reserve()} 的 CPU 本地序号分配机制以精确区分内核缓冲区丢包与用户态丢失，结合 SHA-256 哈希链持久化、系统调用返回值校验的跨智能体文件争用检测（\code{MULTI\_AGENT\_FILE\_CONTENTION}）以及审计优先的 Web 服务派生 Shell 启发式规则；最后，建立\textbf{解耦的三层确定性执行控制与提示词外（Out-of-Prompt）隐私网关}，协同用户态命令策略（\code{ALLOW/BLOCK/ALERT/REWRITE}）、内核 \code{cgroup/connect4,6} 与 \code{sendmsg4,6} 网络阻断及 14 挂载点 BPF LSM 文件/执行阻断，并引入防数据泄漏的行为分组（\code{grouped\_behavior\_v1}）机器学习评估协议。在 100 周期、18 类系统调用共 2,700 万次微基准测量中，系统 eBPF 钩子平均耗时变化为 $-0.95\%$（处于 $-5.26\%\sim +4.05\%$ 测量噪声带内，单次调用绝对增量 $<60\,$ns）；在 19 项涵盖正常开发、恶意利用与智能体特有攻击的运行时回放场景中实现 100\% 上下文继承与精确告警匹配。

\vspace{4pt}
\noindent\textbf{关键词：}AI智能体安全；eBPF；BPF LSM；跨层行为归因；运行时强制控制；隐私网关；数据泄漏审计
\vspace{10pt}
\end{abstract}
\end{@twocolumnfalse}
]

% ============================================================================
\section{引言}
随着大语言模型（LLM）向行动型智能体（Agentic AI）演进，以 Claude Code、Codex CLI、Gemini CLI、Cursor 等为代表的自主编程智能体以及模型上下文协议（MCP）插件已被广泛部署于开发者工作站与持续集成（CI）节点。与传统静态代码补全不同，自主智能体能够连续调用命令执行、文件读写、依赖安装与网页检索工具。一旦智能体遭遇嵌入在代码仓库 Issue、Pull Request 补丁、第三方依赖脚本或网页中的间接提示词注入（Indirect Prompt Injection），原本声明为“只读代码审查”或“读取网页”的工具调用便可能在底层派生异常子进程、窃取 \code{\textasciitilde/.ssh/id\_rsa} 与 \code{.env} 凭证、建立反向 Shell 或通过 MCP 工具隐蔽外发隐私数据\cite{sope2026,flowseal2026,privacypeek2026}。

为实现智能体运行时的安全可信审计与控制，系统安全平面必须针对每一次智能体运行回答四个核心问题：
\begin{enumerate}[label=(\arabic*)]
  \item \textbf{声明意图（Declared Intent）：}智能体在应用层声明调用了何种工具及参数？
  \item \textbf{进程事实（Process Reality）：}操作系统底层实际派生了哪些子进程与可执行文件？
  \item \textbf{资源足迹（Resource Footprint）：}哪些敏感文件路径与外部网络端点被实际打开、修改或连接？
  \item \textbf{控制与证据边界（Control \& Evidence Boundary）：}哪些异常触发了异步语义告警，哪些操作在内核控制点被同步阻断，且证据链是否完整防篡改？
\end{enumerate}

现有安全机制难以同时回答上述问题。一方面，应用层权限弹窗或提示词级隐私指令运行在命令下发之前，无法观测编译脚本、测试框架（如 \code{pytest}、\code{npm install}）内部派生的多层操作系统行为，且模型内提示词极易被上下文清洗（Context Laundering）绕过\cite{flowseal2026}。另一方面，以 Falco\cite{falco}、Cilium Tetragon\cite{tetragon}、Aqua Tracee\cite{tracee} 为代表的云原生 eBPF 安全工具虽能精细捕获系统调用，但缺乏智能体会话、轮次（Turn）与工具调用（Tool Call）的语义上下文；当多个智能体或子智能体在同一主机并发执行同名进程（如 \code{git}、\code{python}、\code{node}）时，仅凭进程名（\code{comm}）或裸 PID 极易发生跨任务误归因。

近期学术界已开始探索面向 AI 智能体的 eBPF 观测技术。AgentSight\cite{agentsight2025} 通过 uprobe 拦截 TLS 明文流量以提取提示词语义并与内核事件关联，但在本地研发环境中常驻开启 TLS 明文捕获会直接暴露私有代码与 API 密钥，带来严重的合规与隐私风险。Agent-Warden\cite{agentwarden2026} 提出基于 eBPF 的内核原生进程—普通文件溯源图（Provenance DAG），侧重于跨文件重命名的事后因果图重建，但未提供跨智能体 CLI 的显式工具调用归因、多模态同步内核阻断以及提示词外隐私网关。

本文在前期工程原型\cite{mingjingrepo}基础上，结合近期完成的 2,700 万次内核开销基准测试、取证级审计链硬化、结果感知语义规则、Web 服务派生 Shell 审计以及防泄漏机器学习语料审计，系统阐述\textbf{明镜高悬（MingJing）}的设计与实证结果。本文主要贡献如下：
\begin{enumerate}
  \item \textbf{三模态显式上下文接入与有界归因方法（第\ref{sec:arch}节）：}支持 Python/Node.js PID 适配器、10 类主流智能体 CLI 原生钩子与 UDS 命令包装器，设计“直接 PID $\to$ 受约束父进程链 $\to$ cgroup v2 $\to$ 未归因保留”的四级归因算法，在默认关闭 TLS 明文劫持下实现细粒度行为核对。
  \item \textbf{取证级遥测与结果感知语义检测（第\ref{sec:telemetry}节）：}提出早于 \code{bpf\_ringbuf\_reserve()} 的 CPU 本地序号机制与代际标识（Generation ID），精确区分内核环形缓冲区满溢与用户态丢包；实现单调时钟对齐、SHA-256 哈希链 JSONL 持久化，以及基于系统调用返回值校验的多智能体文件争用规则与 Web 服务 Shell 派生检测。
  \item \textbf{三层解耦强制控制与提示词外隐私网关（第\ref{sec:control}节）：}将用户态异步语义告警与内核态同步硬阻断严格解耦，协同 \code{agent-wrapper}（\code{ALLOW/BLOCK/ALERT/REWRITE}）、\code{cgroup/connect4,6} + \code{sendmsg4,6} 网络控制与 14 挂载点 BPF LSM 文件/执行控制；构建支持 RFC 6901 JSON 指针策略、重复键拒绝与局部占位符还原的 HTTP 隐私网关。
  \item \textbf{多维实证评估与机器学习数据泄漏审计（第\ref{sec:eval}节）：}完成 100 周期 2,700 万次系统调用开销测试、19 项运行时回放测试与 140,126 组机器学习超参数网格评估；揭示静态命令数据集中的行随机划分泄漏与上游误标注问题，提出基于行为签名分组（\code{grouped\_behavior\_v1}）、保留高危目标路径（\code{/etc/shadow} 等）与导入域局部类型化伪名（\code{\_\_USER\_0001\_\_}）的数据清洗方法。
\end{enumerate}

% ============================================================================
\section{威胁模型与研究定位}
\label{sec:threat}

\subsection{保护资产与威胁来源}
系统旨在保护整个\textbf{智能体执行链（Agent Execution Chain）}：包括工作区源码与构建产物、宿主机敏感凭证（\code{\textasciitilde/.ssh}、\code{.env}、云凭证、\code{/etc/shadow}）、工作区外文件系统边界、智能体派生进程的出站网络连接，以及审计日志本身的防篡改完整性。威胁来源涵盖：（1）本地文件、PR、网页中的间接提示词注入；（2）恶意第三方依赖安装脚本或恶意 MCP 工具服务器；（3）被劫持的本地 Web 服务工作进程（PHP、Python WSGI/ASGI、Node.js）派生交互式 Shell；（4）多智能体并发交错修改共享文件或陷入进程分叉风暴（Fork Storm）；（5）通过 JSON 重复键、双重 URL 编码或历史记忆字段（\code{taskHistory}）实施的隐私外泄。

\subsection{信任边界与显式非目标（Non-Goals）}
本文假设宿主机 Linux 内核（支持 BTF、Tracepoint、cgroup v2 及可选 BPF LSM）与高权限后端守护进程可信；智能体进程、子进程以及钩子自行上报的任务语义均被视为\textbf{待核验的不可信声明}。为保持严谨的学术证据边界，本文明确声明以下\textbf{非目标}：
\begin{itemize}
  \item \textbf{不防御 root 级攻击者或失陷内核：}拥有 root 权限的对手可卸载 BPF 程序或篡改 \code{/sys/fs/bpf/agent-ebpf/} 下的固定映射表（Pinned Maps）。
  \item \textbf{精确映射表不等于完整容器沙箱：}内核 \code{cgroup} 阻断表采用精确 IPv4/IPv6/端口/cgroup inode 键值（不支持 CIDR 网段或 L7 域名解析防火墙），BPF LSM 文件规则当前基于文件/目录 basename（非递归目录树沙箱）。
  \item \textbf{既有状态不回溯切断：}阻断规则下发前已建立的 TCP 连接不会被回溯断开；规则生效前已建立的可读写 \code{mmap} 内存映射不会自动撤销（但后续新的 \code{mmap\_file} 或 \code{file\_mprotect} 权限变更将被拒绝）。
\end{itemize}

% ============================================================================
\section{系统架构与跨层行为归因}
\label{sec:arch}

系统由上下文接入层、eBPF 内核采集层、Go 后端关联与审计层、多层执行控制与隐私网关、以及多端可视化层（Vue 3 Web 工作台、原生 Renew 桌面端、终端 TUI 及 MCP/OTLP 接口）组成。

\subsection{三模态语义上下文接入}
为在不开启 TLS 明文解密的前提下获取任务边界，系统设计三类互补接入通道：
\begin{enumerate}
  \item \textbf{语言运行时 PID 适配器：}提供 Python（\code{agent\_tracker.py}）与 Node.js（\code{agentTracker.js}）轻量级 SDK，在进程启动与退出时通过 \code{POST /register} 与 \code{/unregister} 维护内核 \code{agent\_pids} 映射表与用户态 \code{ProcessContextStore}。
  \item \textbf{10 类主流智能体 CLI 原生生命周期钩子：}原生适配 Claude Code、Codex、Gemini CLI、GitHub Copilot、Kiro CLI、Pi、Oh My Pi、Augment、Antigravity CLI 与 ZCode，在工具调用前后注入 \code{agent\_run\_id}、\code{task\_id}、\code{tool\_call\_id}、\code{trace\_id}、工作目录 \code{cwd} 与参数摘要 \code{argv\_digest}。
  \item \textbf{基于对端凭证校验的 UDS 命令包装器（\code{agent-wrapper}）：}针对 Cursor、DeepSeek Harness（\code{dsh}）等包装场景，在命令执行入口通过私有 Unix 域套接字 \code{/tmp/agent-ebpf.sock} 发送 \code{pb.WrapperRequest}；后端 \code{server\_uds.go} 校验本地套接字 peer credentials 后执行同步策略决策。
\end{enumerate}

\subsection{有界上下文关联算法与进程图谱边界}
内核采集程序（\code{agent\_tracker.c}）挂载于 \code{execve}、\code{openat}、\code{connect}、\code{mkdirat}、\code{unlinkat}、\code{ioctl}、\code{bind}、\code{sendto}、\code{recvfrom} 等 9 类系统调用 tracepoint，通过零拷贝 BPF Ring Buffer 上报事件。后端 \code{events.EnrichEventContext} 按四级优先级执行归因：
\[
\footnotesize
\begin{aligned}
\text{当前 PID 上下文} &\to \text{受创建时间约束的可信父进程上下文} \\
&\to \text{cgroup v2 归属} \to \text{保留为未归因}.
\end{aligned}
\]
系统严格遵循两项归因红线：（1）相同 \code{comm} 名称或普通 shell 执行绝不自动推定为某智能体身份；（2）桌面端（\code{native\_agent\_lineage\_model.go}）明确区分\textbf{操作系统进程快照}与\textbf{智能体因果图}，利用进程启动时间防范 PID 复用污染，并将启发式派生节点显式标注为 \code{subagent-candidate}。

% ============================================================================
\section{取证级遥测与结果感知语义检测}
\label{sec:telemetry}

\subsection{Pre-Reserve 丢包溯源、双时钟对齐与哈希链持久化}
\paragraph{1. 早于 Ringbuf 预留的序号分配。}
传统 eBPF 采集器仅在 \code{bpf\_ringbuf\_reserve()} 成功后分配序号，导致缓冲区满溢只能通过宏观计数器粗略估计。本系统在调用 \code{bpf\_ringbuf\_reserve()} \textbf{之前}即分配 CPU 本地单调递增序号 \code{kernel\_sequence}。若预留失败，该序号被消耗并累加每 CPU 的 \code{pending\_dropped\_events}；下一条成功事件将携带重载代际标识 \code{kernel\_audit\_generation}、距上次成功事件的丢包数 \code{kernel\_dropped\_since\_last} 及累计预留失败数 \code{kernel\_reserve\_failures\_total}。借此，后端可通过 $(\code{generation}, \code{cpu}, \code{sequence})$ 严格区分\textbf{内核环形缓冲区满溢丢包}与\textbf{用户态读取/解码丢失}，并通过 C \code{\_Static\_assert}、Go \code{unsafe.Offsetof} 测试及 CI 内核验证器三重校验保障 ABI 一致。

\paragraph{2. 双时钟对齐与零分配 TLS 分片重组。}
系统严格区分内核单调时钟捕获时刻（\code{probe\_timestamp\_ns}，来自 \code{bpf\_ktime\_get\_ns()}）与用户态墙上时钟接收时刻（\code{ingest\_timestamp\_ns}），计算内核至用户态处理延迟 \code{capture\_delay\_ns}，并在离线回放时显式标记 \code{unix\_replay}。在可选 TLS uprobe 诊断链路（支持 OpenSSL、GnuTLS、NSS、Go \code{crypto/tls}、Rustls 五大协议栈）中，用户态重组器为每条待处理消息仅分配一次固定槽位缓冲区，利用 32 位位图（Bitset）追踪最多 18 个分片并原地紧凑化（In-place Compaction），消除了高并发流式输出下的分片级 Map 内存分配与 GC 压力。

\paragraph{3. 防篡改 SHA-256 哈希链持久化。}
JSONL 事件持久化模块为每个写入代际（Writer Generation）建立随机链标识 \code{auditChainId}，每条记录的 \code{auditHash} 由链元数据、接收时间戳、前一条记录哈希 \code{auditPrevHash} 以及 \code{EventEnvelope} 的确定性 Protobuf 编码经 SHA-256 计算生成，可在不引入二次 JSON 序列化开销的前提下通过 \code{VerifyAuditChain} 精确检测记录篡改与删减。

\subsection{结果感知的语义安全告警规则}
用户态规则引擎（\code{alerts\_semantic.go}）维护有界历史状态并合成 10 类告警（表\ref{tab:alerts}）。针对传统规则易产生误报的痛点，本文重点实现两项高精度机制：

\begin{table}[t]
\centering
\footnotesize
\caption{用户态合成的 10 类语义安全告警及判定焦点（均与内核同步阻断解耦）。}
\label{tab:alerts}
\begin{tabularx}{\columnwidth}{@{}lX@{}}
\toprule
\textbf{告警代码} & \textbf{核心触发语义与证据要求} \\
\midrule
\code{SECRET\_ACCESS} & 访问或尝试读取 \code{.ssh}、\code{.env}、\code{/etc/shadow} 等敏感凭证路径 \\
\code{SEMANTIC\_MISMATCH} & 声明的工具意图（如只读检查）与实际 OS 系统调用行为冲突 \\
\code{UNEXPECTED\_NETWORK\_EGRESS} & 未声明网络需求的工具发生外部网络连接 \\
\code{UNEXPECTED\_CHILD\_PROCESS} & 派生未在信任链注册的异常子进程（如 \code{curl}、\code{ssh}） \\
\code{TOOL\_BEHAVIOR\_DRIFT} & 工具执行序列偏离基线模式（如 \code{chmod +x} 后立即 \code{execve}） \\
\code{SUSPICIOUS\_SHELL\_PIPELINE} & 含危险管道或重定向的 Shell 命令（如 \code{curl | sh}、\code{/dev/tcp}） \\
\code{WORKSPACE\_ESCAPE} & 越权读写当前工作空间（\code{cwd}）以外的系统目录 \\
\code{TOKEN\_EXFIL\_RISK} & 检测到读取敏感文件后紧接向外部端点发起疑似外发连接 \\
\code{RESOURCE\_WASTING\_LOOP} & 短周期内进程分叉风暴（Fork Storm）或重复外部请求循环 \\
\code{MULTI\_AGENT\_FILE\_CONTENTION} & 两个不同身份 Agent 在 15\,s 内对同一容器域绝对路径成功交错写入 \\
\bottomrule
\end{tabularx}
\end{table}

\begin{enumerate}
  \item \textbf{基于返回结果校验的跨智能体文件争用（\code{MULTI\_AGENT\_FILE\_CONTENTION}）：}仅当（i）写操作或元数据修改系统调用\textbf{实际成功}（\code{retval >= 0}，且 \code{write(2)} 必须返回正字节数 \code{retval > 0}，拒绝失败或零字节写入）；（ii）路径为可解析的真实普通文件绝对路径（排除 \code{fd:3} 占位符、套接字及 \code{/dev/null}、\code{/proc/*}、\code{/sys/*}、\code{/dev/pts/*} 等伪文件，保留 \code{/dev/shm} 普通文件）；（iii）两次写入具备可比较且不同的根智能体身份（不同 \code{root\_agent\_pid} 或 \code{agent\_run\_id}）；（iv）按 $(\code{container\_id}, \code{path})$ 隔离容器作用域且忽略乱序旧时间戳；（v）处于 15\,s 时间窗与 15\,s 冷却期内时才触发告警（风险下限设为 $70/100$，明确定义为路径级交错修改线索而非内核数据竞争证明）。
  \item \textbf{审计优先的 Web 服务派生 Shell 检测（\code{web\_shell.go}）：}通过有界 procfs 查询父进程命令行，仅当父进程具备明确的 PHP（\code{php-fpm}、\code{artisan serve}）、Python（\code{uvicorn}、\code{gunicorn}、Flask/FastAPI/Django \code{runserver}）或 JS/TS（Node/Bun/Deno 服务入口、Next/Nuxt/Nest）Web 服务证据且子进程精确匹配 Shell 名称（避免将 \code{shellcheck} 误判为 \code{sh}）时触发，从而使 \code{python -m pytest} 或 \code{node scripts/build.js} 派生 Shell 保持零误报。
\end{enumerate}

% ============================================================================
\section{多层确定性执行控制与提示词外隐私网关}
\label{sec:control}

\subsection{三层解耦阻断与受控内核风险反馈}
系统构建三层确定性控制面：
\begin{enumerate}
  \item \textbf{用户态 Wrapper 策略：}在 \code{agent-wrapper} 入口对精确 \code{comm} 与正则表达式执行 \code{ALLOW}、\code{BLOCK}（退出码 1）、\code{ALERT} 或命令行参数重写 \code{REWRITE}。
  \item \textbf{内核态 cgroup v2 网络硬阻断（\code{cgroup\_sandbox.c}）：}挂载于 \code{cgroup/connect4}、\code{connect6}、\code{sendmsg4}、\code{sendmsg6}，基于 cgroup v2 inode ID、精确 IPv4（含 \code{::ffff:a.b.c.d} 归一化）、IPv6 及 TCP/UDP 端口映射表，在内核态同步返回 \code{-EPERM}，覆盖未连接 UDP \code{sendto/sendmsg} 与已连接 UDP 发送。
  \item \textbf{内核态 14 挂载点 BPF LSM 强制访问控制（\code{lsm\_enforcer.c}）：}覆盖执行阶段 \code{bprm\_check\_security}（精确路径或可执行文件 basename）及 13 个文件与 inode 钩子（\code{file\_open}、\code{file\_permission}、\code{mmap\_file}、\code{file\_mprotect}、\code{inode\_setattr}、\code{inode\_create}、\code{inode\_link}、\code{inode\_symlink}、\code{inode\_unlink}、\code{inode\_mkdir}、\code{inode\_rmdir}、\code{inode\_mknod}、\code{inode\_rename}），对命中目标同步返回 \code{EACCES}，有效拦截通过已打开文件描述符执行的读写、\code{ftruncate}/\code{fchmod} 及内存映射权限提升。
\end{enumerate}

为防止高风险评分误锁死宿主机，系统规定风险评分与策略提案（\code{agentboundary.PlanTransition}）默认不修改 BPF Map；异步反馈协程（\code{KernelRiskFeedbackWorker}）仅在 \code{PolicyManagementEnabled} 与 \code{KernelRiskFeedback.Enabled} 双重运行时门控开启且分值达到阈值时才向内核下发动作。

\subsection{提示词外 HTTP 隐私网关与防泄漏机器学习清洗}
针对近期文献揭示的智能体获取阶段隐私泄漏与 MCP 隐蔽外传风险\cite{flowseal2026,minim2026,sope2026,privacypeek2026}，系统在四级日志脱敏（\code{None/Basic/Standard/Strict}）之外实现了运行于 LLM 上下文外部的 HTTP 隐私网关：（1）采用有界深度/节点数的 JSON 解析器并拒绝任何包含重复键（含 Unicode 等价键名）的请求；（2）严格模式下默认拒绝 \code{taskHistory} 等全量记忆导出字段；（3）支持精确到 RFC 6901 JSON 指针的 \code{block}、\code{drop}、\code{redact}（生成请求域内不可预测占位符并在本地 HTTP/SSE 响应中按需还原）及绑定真实解析上游主机名的 \code{allow} 规则；（4）银行卡检测强制要求合法发卡行前缀与 Luhn 校验和通过，避免将 12--15 位普通数字 ID 误判为敏感信息。

在机器学习训练管线（\code{decision\_training.go}）中，本文识别并修复了三项影响学术可信度的数据泄漏缺陷：
\begin{enumerate}
  \item \textbf{行为签名分组留出（\code{grouped\_behavior\_v1}）：}使用 \code{PrepareAutoTuneSplit} 按规范化命令与行为特征签名分组，确保同模板重复样本绝不跨越训练集与验证集。
  \item \textbf{保留风险目标路径的脱敏（\code{dataset\_sanitize.go}）：}在掩码真实密钥、Bearer 令牌、密码哈希与私有 IP 的同时，显式保留 \code{/etc/passwd}、\code{/etc/shadow}、\code{/etc/sudoers}、\code{/etc/ssh/sshd\_config} 等公开风险目标路径名及相对路径穿越特征（如 \code{../../etc/shadow}），避免粗暴路径抹除导致模型丧失攻击目标分辨力。
  \item \textbf{导入域局部类型化伪名（\code{dataset\_identity.go}）：}将同一导入批次内的用户名与项目名映射为一致的类型化符号（如 \code{/home/\_\_USER\_0001\_\_/projects/\_\_PROJECT\_0001\_\_}），既保留跨参数结构关联，又杜绝跨数据集全局哈希碰撞关联。
\end{enumerate}

% ============================================================================
\section{实验评估与结果分析}
\label{sec:eval}

\subsection{eBPF 内核钩子微基准开销实验（2,700 万次测量）}
在搭载 Linux \code{7.0.12-1-cachyos} 内核、已挂载 186 个 tracepoint 的现代 x86\_64 平台上，我们对 5 大类共 18 个系统调用执行了 \textbf{100 个完整周期}的基线与 eBPF 交替对照测试（每周期每项操作预热 1 次共 5,000 次调用，正式测量 3 轮×5,000 次调用，总计 \textbf{2,700 万次系统调用测量}，采用 Welford 在线算法聚合）。

\begin{table}[t]
\centering
\scriptsize
\setlength{\tabcolsep}{2.8pt}
\caption{100 周期、2,700 万次测量下 eBPF 内核钩子耗时对照（Linux 7.0.12-1-cachyos）。}
\label{tab:overhead_zh}
\begin{tabular}{@{}llrrrr@{}}
\toprule
\textbf{类别} & \textbf{系统调用} & \textbf{基线($\mu$s)} & \textbf{eBPF($\mu$s)} & \textbf{变化(\%)} & \textbf{CV(\%)} \\
\midrule
\multirow{7}{*}{进程}
 & \code{getpid}       & 0.141 & 0.137 & $-3.05$ & 34.1 \\
 & \code{getppid}      & 0.121 & 0.115 & $-4.69$ & 34.7 \\
 & \code{gettid}       & 0.107 & 0.104 & $-3.00$ & 35.4 \\
 & \code{getuid}       & 0.099 & 0.097 & $-2.29$ & 33.7 \\
 & \code{getgid}       & 0.099 & 0.096 & $-3.67$ & 33.4 \\
 & \code{geteuid}      & 0.098 & 0.098 & $-0.02$ & 34.1 \\
 & \code{getegid}      & 0.093 & 0.097 & $+4.05$ & 29.8 \\
\midrule
\multirow{6}{*}{文件}
 & \code{open\_close}  & 1.063 & 1.049 & $-1.34$ & 14.8 \\
 & \code{stat}         & 0.620 & 0.625 & $+0.82$ & 16.6 \\
 & \code{access}       & 0.584 & 0.591 & $+1.31$ & 16.5 \\
 & \code{getcwd}       & 1.629 & 1.609 & $-1.24$ & 13.3 \\
 & \code{read\_write}  & 5.757 & 5.787 & $+0.53$ & 14.6 \\
 & \code{readlink}     & 1.045 & 0.990 & $-5.26$ & 14.4 \\
\midrule
\multirow{2}{*}{网络}
 & \code{socket\_close}& 1.468 & 1.430 & $-2.60$ & 13.3 \\
 & \code{bind}         & 2.541 & 2.523 & $-0.70$ & 13.9 \\
\midrule
目录 & \code{mkdir\_rmdir} & 3.265 & 3.322 & $+1.73$ & 10.5 \\
\midrule
\multirow{2}{*}{控制}
 & \code{clone}        & 0.100 & 0.103 & $+3.29$ & 33.7 \\
 & \code{prctl}        & 0.270 & 0.267 & $-0.96$ & 19.6 \\
\midrule
\multicolumn{2}{@{}l}{\textbf{18 项总体平均}} & --- & --- & $\mathbf{-0.95}$ & \textbf{23.7} \\
\bottomrule
\end{tabular}
\end{table}

如表\ref{tab:overhead_zh}所示，18 项系统调用的平均相对耗时差为 $-0.95\%$，分布在 $-5.26\%$（\code{readlink}）至 $+4.05\%$（\code{getegid}）之间，没有任何单项操作开销超过 $+5\%$，单次调用的最大绝对耗时增量仅为 $0.057\,\mu$s（\code{mkdir\_rmdir}）。从系统测量方法学角度分析，对于耗时仅 $\sim 0.1\,\mu$s 的极快调用（CV 约 $30\%\sim 35\%$），负向百分比（$-0.003\sim -0.055\,\mu$s）并非 eBPF 真正“加速”了内核，而是表明 eBPF 过滤逻辑的开销已完全落入 CPU 缓存局部性与系统调度噪声带之内。相比单周期测试（平均 $+10.68\%$，波动 $-16\%\sim +86\%$），100 周期交替测量成功消除了系统热漂移，证明本系统内核采集层在常规研发负载下开销可忽略。

\subsection{19 项运行时回放场景与多模型扫描评估}
\paragraph{1. 智能体运行时回放基准（\code{runtime-replay.v1}）。}
我们对 \code{benchmarks/runtime-replay/scenarios.json} 中的 19 个标准场景进行了自动化回归测试：
\begin{itemize}
  \item \textbf{6 类正常研发工作流}（\code{git status}、\code{npm install}、\code{pip install}、\code{pytest}、\code{cargo build}、\code{rg} 只读代码审查）：子进程 100\% 正确继承 \code{agent\_run\_id} 与 \code{tool\_call\_id}，Wrapper 100\% 判定为 \code{ALLOW}，且实现\textbf{零误报}。
  \item \textbf{8 类典型命令/系统利用}（\code{curl | sh} 管道、读取 \code{id\_rsa}、\code{/dev/tcp} 反向 Shell、访问 \code{/etc/shadow} 逃逸、\code{chmod +x} 执行临时载荷、异常 \code{ssh} 横向移动、隐蔽 Webhook 外连、9 子进程 Fork 风暴）：100\% 命中预期告警组合。
  \item \textbf{5 类智能体特有攻击}（网页提示词注入导致读取 \code{.env} 并外发至 Pastebin、恶意 MCP 插件调用 \code{curl}、前端 DOM 检查开启隐蔽遥测外连、组件编辑任务越权跳转远程 SSH 开发机、重构任务陷入进程派生死循环）：100\% 完成跨进程归因并触发全部目标告警码（如 \code{TOKEN\_EXFIL\_RISK}、\code{RESOURCE\_WASTING\_LOOP}）。
\end{itemize}

\paragraph{2. 140,126 组机器学习模型扫描与外部语料审计。}
在本地 949 条标注命令集与 93 条均衡集上，我们对 42 个内置模型配置进行了单参数 1,000 点全面网格扫描（共 190 项覆盖清单、140,126 组训练结果）。在可比留出评估（Holdout-comparable）下，轻量随机森林（\code{random\_forest\_fast}，5 棵树，深度 4）在 0.064\,s 训练耗时下达到 $7.97\times 10^6$ 次/秒的推理吞吐（单样本 $<0.3\,\mu$s），且在 10 次重复稳定性复测中 \code{svm\_long}（$\text{lr}=0.1, \text{iter}=8000$）与随机森林均保持高验证准确率与 100\% 的 \code{ALLOW} 放行率；而未经调优的 \code{extra\_trees} 与默认感知机虽具备一定整体准确率，但其 \code{ALLOW} 放行率跌至 $0\%$，说明在类别偏斜的安全策略场景中必须将良性请求放行率作为核心一等指标。

同时，我们对公开 \code{bash-datasets}（\code{taglio4\_3k}）的 124 条样本开展了逐命令审计，发现上游数据集存在 3 处明显反向误标（如将 \code{find | xargs rm -rdf} 和 \code{shred} 标为安全，将收紧私钥权限的 \code{chmod 0600 id\_rsa} 标为危险），且有 25 条命令（20.2\%）属于必须依赖运行时工作目录与权限上下文才能裁决的 \code{DEFER} 类别。这一发现从实证角度证明：\textbf{静态命令风险模型不能单独作为阻断门控，必须与本文提出的内核跨层归因及确定性 OS 策略层协同工作。}

\subsection{公开基准评测与核心组件消融实验}
为进一步验证系统在大规模外部公开语料上的泛化能力并定量剥离各核心组件的独立贡献，我们引入学术界公认的三大公开基准——\textbf{InjecAgent}（1,054 项间接提示词注入用例）\cite{injecagent2024}、\textbf{PrivacyLens}（493 条智能体隐私轨迹 + 493 项良性对照）\cite{privacylens2024} 及全量去重 \textbf{bash-datasets}（2,187 条命令）——联合多智能体并发归因、文件争用与 Web-Shell 规则完成系统化消融（测试套件 \code{TestExternalPublicBenchmarksAndAblations}）。

\paragraph{1. 三模态上下文接入与跨层归因消融（88 条子进程/并发多跳事件）。}
在 19 项 \code{runtime-replay.v1} 场景与 16 个并发多智能体 3 跳进程树（共享 \code{python}/\code{git}/\code{node}/\code{bash}）上，表\ref{tab:ablation_attr_zh} 分别报告 \textbf{Run 级归因率} 与 \textbf{Tool-Call 级归因率}：朴素 \code{comm} 启发式产生 \textbf{13.6\%} 跨任务误归因（并发叶节点误归因率达 75\%）；仅 PID 接入虽恢复 100\% Run 级归属，但工具级归因为 0\% 且告警召回率仅 63.6\%；而\textbf{完整明镜三模态联合归因}在两级归因、告警召回与良性零误报上均达到 \textbf{100.0\%}。

\begin{table}[t]
\centering
\scriptsize
\setlength{\tabcolsep}{2.0pt}
\caption{消融一：5 种接入模式在 88 条子进程/16 并发多跳事件上的双层归因与告警对比。}
\label{tab:ablation_attr_zh}
\begin{tabular}{@{}lrrrrr@{}}
\toprule
\textbf{接入模式} & \textbf{Run率} & \textbf{Tool率} & \textbf{误归因} & \textbf{告警召回} & \textbf{FP} \\
\midrule
无上下文（仅内核）       & 0.0\%   & 0.0\%   & 0.0\%  & 57.6\% & 0 \\
朴素 \code{comm} 启发式  & 76.1\%  & 0.0\%   & 13.6\% & 60.6\% & 0 \\
仅 PID（无工具意图）     & 100.0\% & 0.0\%   & 0.0\%  & 63.6\% & 0 \\
仅 Hook/Wrapper（无多跳）& 0.0\%   & 0.0\%   & 0.0\%  & 57.6\% & 0 \\
\textbf{完整明镜三模态}  & \textbf{100.0\%} & \textbf{100.0\%} & \textbf{0.0\%} & \textbf{100.0\%} & \textbf{0} \\
\bottomrule
\end{tabular}
\end{table}

\paragraph{2. InjecAgent 公开基准评测（1,054 项间接提示词注入）。}
在全部 1,054 项 \textbf{InjecAgent} 用例（510 项直接危害 DH + 544 项数据窃取 DS）上：仅靠内核事件（无工具声明意图）能通过敏感读取-外连关联捕获 100\% 数据窃取，但对直接危害（DH）的检出率为 \textbf{0.0\%}（总检出率 51.6\%）；引入明镜跨层工具意图关联后，DH 检出率跃升至 \textbf{94.1\%}（480/510），DS 保持 \textbf{100.0\%}（544/544），总检出率达 \textbf{97.2\%}（1,024/1,054），单用例 P50/P95 判定延迟仅为 \textbf{5.39\,$\mu$s / 11.74\,$\mu$s}。

\paragraph{3. 结果感知文件争用与 Web-Shell 启发式消融。}
表\ref{tab:ablation_sem_zh} 展示两项语义规则的硬化收益：（i）在 30 组文件争用/噪声集上，明镜全部结果感知护栏将误报从 20/20 降至 \textbf{0}（$F_1$ 从 $0.500$ 升至 \textbf{1.000}）；（ii）在 31 组真实运行时进程对（17 组 PHP/Python/Node Web 服务派生 Shell + 14 组 \code{pytest}/\code{npm run dev}/\code{shellcheck} 等良性 CLI）上，明镜正向 Web 服务证据规则将误报从 14/14 降至 \textbf{0}（精确率由 54.8\% 升至 \textbf{100.0\%}，$F_1$ 由 0.708 升至 \textbf{1.000}）。

\begin{table}[t]
\centering
\scriptsize
\setlength{\tabcolsep}{2.4pt}
\caption{消融二：多智能体文件争用（30 组）与 Web-Shell 启发式规则（31 组）消融对比。}
\label{tab:ablation_sem_zh}
\begin{tabular}{@{}lrrrrr@{}}
\toprule
\textbf{检测器变体} & \textbf{TP} & \textbf{FP} & \textbf{精确率} & \textbf{召回率} & \textbf{F1} \\
\midrule
文件争用：朴素 Path+PID     & 10/10 & 20/20 & 33.3\% & 100.0\% & 0.500 \\
文件争用：移除容器/伪文件   & 10/10 & 10/20 & 50.0\% & 100.0\% & 0.667 \\
文件争用：移除返回值校验    & 10/10 & 5/20  & 66.7\% & 100.0\% & 0.800 \\
\textbf{文件争用：完整明镜} & \textbf{10/10} & \textbf{0/20} & \textbf{100.0\%} & \textbf{100.0\%} & \textbf{1.000} \\
\midrule
Web-Shell：朴素解释器规则   & 17/17 & 14/14 & 54.8\% & 100.0\% & 0.708 \\
\textbf{Web-Shell：完整明镜}& \textbf{17/17} & \textbf{0/14} & \textbf{100.0\%} & \textbf{100.0\%} & \textbf{1.000} \\
\bottomrule
\end{tabular}
\end{table}

\paragraph{4. 全量 2,187 条 Bash 语料防泄漏划分、脱敏与均衡采样消融。}
在全量去重 2,187 条 \textbf{bash-datasets}（10 种子，72 棵树深度 14）与 48 种子 1:49 稀有攻击集上（表\ref{tab:ablation_ml_zh}）：（i）在 2 倍重复日志（4,374 行）上，行随机划分导致 \textbf{688 对命令跨集泄漏}，使验证准确率虚高至 $98.63\% \pm 0.33\%$；明镜 \code{grouped\_behavior\_v1} 实现 \textbf{0 泄漏}，真实验证准确率为 $97.76\% \pm 0.39\%$，BLOCK 召回率 $91.25\% \pm 2.14\%$（Wilson 95\% 置信区间 $[80.7\%, 96.1\%]$），良性 FPR 仅 $1.29\% \pm 0.38\%$。（ii）若粗暴脱敏抹除 \code{/etc/shadow} 等高危路径，BLOCK 召回率暴跌 \textbf{23.04\%}（$91.25\% \to 68.21\% \pm 3.01\%$），FPR 翻倍以上（$1.29\% \to 3.02\%$）。（iii）在 1:49 稀有攻击集上，类别分层 Bootstrap 将召回率由 $52.08\%$ 提升至 \textbf{100.0\%}（48/48，0 误报）。

\begin{table}[t]
\centering
\scriptsize
\setlength{\tabcolsep}{1.8pt}
\caption{消融三：全量 2,187 条 Bash 语料（10 种子均值$\pm$标准差）与稀有攻击采样（48 种子）。}
\label{tab:ablation_ml_zh}
\begin{tabular}{@{}lrrrrr@{}}
\toprule
\textbf{ML 管线变体} & \textbf{泄漏} & \textbf{验证率(\%)} & \textbf{阻断召回(\%)} & \textbf{FPR(\%)} & \textbf{均衡(\%)} \\
\midrule
朴素行划分 (2x, 4.3k)   & 688 & 98.63$\pm$0.33 & 92.63$\pm$1.18 & 0.47$\pm$0.28 & 96.08 \\
\textbf{明镜分组 (2,187)}& \textbf{0} & \textbf{97.76$\pm$0.39} & \textbf{91.25$\pm$2.14} & \textbf{1.29$\pm$0.38} & \textbf{94.98} \\
移除 \code{/etc} 保留   & 0   & 93.30$\pm$0.67 & 68.21$\pm$3.01 & 3.02$\pm$0.51 & 82.60 \\
\midrule
均匀采样 (1:49, 48种子) & 0   & ---            & 52.08          & 0.00          & 76.04 \\
\textbf{均衡采样 (48种子)}& \textbf{0} & ---   & \textbf{100.0} & \textbf{0.00} & \textbf{100.0} \\
\bottomrule
\end{tabular}
\end{table}

\paragraph{5. PrivacyLens 公开隐私轨迹评测（493 敏感 + 493 良性对照）。}
在全部 493 条 \textbf{PrivacyLens} 智能体隐私轨迹与 493 条同结构良性对照请求（共 986 次请求，经 50 次预热消除冷启动后测稳态延迟）上：关闭模式与仅审计模式外泄率均为 100.0\%（493/493）；\textbf{明镜脱敏+局部还原模式}将上游外泄率降至 \textbf{0.0\%}（0/493），100.0\%（493/493）通过本地往返还原校验且保持 100.0\% 良性业务可用性，稳态 P50/P95 延迟仅为 \textbf{10.95\,$\mu$s / 16.80\,$\mu$s}；\textbf{严格阻断模式}外泄率 \textbf{0.0\%}、良性可用性 \textbf{100.0\%}，P50/P95 延迟低至 \textbf{8.77\,$\mu$s / 10.23\,$\mu$s}。

% ============================================================================
\section{局限性与结论}
\label{sec:conclusion}

\paragraph{局限性说明。}
本文严格区分已验证能力与边界：（1）当前 \code{MULTI\_AGENT\_FILE\_CONTENTION} 为带容器域与返回值约束的路径级时序相关性检测，尚未实现 Agent-Warden\cite{agentwarden2026} 式的内核 inode 局部存储（BPF Local Storage）污点传播图；（2）内核 \code{cgroup} 与 BPF LSM 规则采用精确键与 basename 匹配，不替代递归目录容器沙箱或 L7 防火墙，也不回溯切断既有 TCP 连接与既有 \code{mmap} 映射；（3）跨多主机、大规模真实开发者会话的长期误报率统计与用户定位耗时对照实验仍作为后续工作推进。

\paragraph{结论。}
本文提出了面向自主编程智能体的 eBPF 跨层行为归因与多层安全控制系统“明镜高悬”。通过三模态显式上下文接入、早于缓冲区预留的取证级序号链、结果感知的语义告警规则、三层解耦的内核/用户态控制以及提示词外 HTTP 隐私网关，系统在保持微秒级以下内核追踪开销的同时，实现了对智能体复杂执行链的透明核对与可靠约束。

% ============================================================================
\footnotesize
\begin{thebibliography}{10}

\bibitem{bpflsm}
Linux Kernel Community.
\newblock LSM BPF Programs[EB/OL].
\newblock \url{https://docs.kernel.org/bpf/prog_lsm.html}, 2026.

\bibitem{falco}
The Falco Authors.
\newblock The Falco Project Documentation[EB/OL].
\newblock \url{https://falco.org/docs/}, 2026.

\bibitem{tetragon}
Cilium Authors.
\newblock Tetragon: eBPF-based Security Observability and Runtime Enforcement[EB/OL].
\newblock \url{https://tetragon.io/docs/}, 2026.

\bibitem{tracee}
Aqua Security.
\newblock Tracee: Runtime Security and Forensics using eBPF[EB/OL].
\newblock \url{https://aquasecurity.github.io/tracee/}, 2026.

\bibitem{agentsight2025}
ZHENG Y, HU Y, YU T, QUINN A.
\newblock AgentSight: System-Level Observability for AI Agents Using eBPF[EB/OL].
\newblock arXiv:2508.02736v2, 2025.

\bibitem{agentwarden2026}
CUI D, GU Z, ZHENG P, HAN S, LIAO Y.
\newblock Agent-Warden: eBPF-Based Kernel-Native Process-File Provenance Tracking for LLM Agents[EB/OL].
\newblock arXiv:2609.38245, 2026.

\bibitem{claudesandbox2025}
DWORKEN D, WELLER-DAVIES O.
\newblock Beyond permission prompts: making Claude Code more secure and autonomous[EB/OL].
\newblock Anthropic Engineering Blog, 2025.

\bibitem{flowseal2026}
FLOWSEAL Authors.
\newblock FLOWSEAL: Confuse the Model, Control the Flow[EB/OL].
\newblock arXiv:2609.14003, 2026.

\bibitem{minim2026}
YU Y, et~al.
\newblock MINIM: Privacy-Aware Minimal View for Agents via Trusted Local Sanitization[C]//Proc. of ICML 2026 (PMLR vol. 306), 2026.

\bibitem{sope2026}
LIN H, et~al.
\newblock SOPE: Situation-Aware and Statistically Indistinguishable Privacy Exfiltration for MCP-enabled Agents[C]//Proc. of ICML 2026 (PMLR vol. 306), 2026.

\bibitem{privacypeek2026}
PrivacyPeek Authors.
\newblock PrivacyPeek: Auditing What LLM-Based Agents Acquire, Not Just What They Say[EB/OL].
\newblock arXiv:2606.00152, 2026.

\bibitem{mingjingrepo}
Agent eBPF Filter 项目仓库.
\newblock 事件协议、策略语义、取证审计与基准测试报告[Z].
\newblock Commit \code{26c28fea}, 2026-11.

\bibitem{injecagent2024}
ZHAN Q, LIANG Z, YING Z, KANG D.
\newblock InjecAgent: Benchmarking Indirect Prompt Injections in Tool-Integrated Large Language Model Agents[C]//Findings of ACL 2024, 2024.

\bibitem{privacylens2024}
SHAO Y, LI T, SHI W, LIU Y, YANG D.
\newblock PrivacyLens: Evaluating Privacy Norm Awareness of Language Models in Action[C]//Advances in Neural Information Processing Systems (NeurIPS 2024), 2024.

\end{thebibliography}

\end{document}
